\section{Discussion, Limitations, and Conclusion}
\label{sec:discussion}

\subsection{Discussion}
Our results show that network robustness depends on what we want the network to preserve. Random edge removal left a large connected core in the 5,000-node BFS graph: at $f=0.50$, the largest connected component still contained 98.2\% of the original nodes, close to the configuration-model reference. However, only 43.3\% of the fixed queries retained a valid two-hop path on the same ten damage states (Section~\ref{sec:functional}). Protein--protein interactions can preserve global connectivity even when a drug and disease lose their last shared-protein route. Thus, connectivity alone is not a sufficient measure of functional robustness.

The type and location of damage also matter. In Section~\ref{sec:structural}, targeted removal of high-degree or high-betweenness nodes reduced the giant component more rapidly than random node removal. At $f=0.50$, the remaining giant-component fractions were approximately 0.232 and 0.206 under these fixed targeted rankings, versus 0.493 under random removal. The two centrality rankings overlapped substantially, but the observed patterns apply to the fixed rankings used here rather than adaptive attacks.

Functional performance also depends on the permitted query horizon. At 50\% random edge removal, the matched three-seed comparison increased retrieval success from 40.3\% at $L=2$ to 84.5\% at $L=3$ (79.2\% over 600 fresh damage orders; Section~\ref{sec:prediction-uncertainty}). Queries with multiple original shared-protein paths were more likely to retain a two-hop route than single-path queries (60.7\% versus 25.9\% at the same damage level). Both path redundancy and allowing additional hops can therefore buffer query failure, although they describe different mechanisms.

Restoration results extend this distinction from robustness to recovery. No single policy performed best for all targets and budgets (Section~7). Edge-betweenness repair quickly restored global connectivity and three-hop reachability; relation-type repair achieved the largest two-hop gain at a 20\% budget after 50\% edge removal (48.50 percentage points). Static task-aware repair was competitive at the smallest budget, but its fixed scores quickly saturated, leaving most of the remaining budget to zero-score ties. These differences show why repair priorities should follow the intended retrieval task and available budget rather than topology alone.

\paragraph{What the theory adds.}
Section~\ref{sec:theory} explains why these outcomes differ, not only that they
differ. Membership of the giant component is a collective property. It depends
on an unbounded number of edges and is described by a self-consistent
generating-function equation with a threshold. For PrimeKG the threshold lies
beyond 99\% edge loss, which is why $S$ hardly moves. A bounded-hop query is a
local property. Its success probability is a finite reliability polynomial,
$R_2=1-G_m(1-p^2)$, with no threshold. Because the two observables are
different kinds of object, a high $S$ cannot certify high $R_2$. Both limits of
the horizon have the same form: one minus the generating function of a
redundancy distribution. At $L=2$ the redundant unit is a two-edge route; at
$L=\infty$ it is a single edge of the minimum cut, which for every query equals
the smaller endpoint degree. Longer horizons therefore move the bottleneck from
shared proteins to the drugs' own target links rather than to the PPI network.
We had expected a percolation-like collapse at long range and did not find one.

These formulas contain no fitted parameters, and they are quantitatively
useful. Under blind damage they match the long-run mean of the simulations
within 0.001 at $L=2$ and within 0.009 at $L=\infty$. They predict E's
independent restoration gains within one seed SD in all twelve cases, and per
query they already reach the AUC ceiling set by chance. They fail where their
one assumption fails. Targeted attacks remove the hubs that the supposedly
independent routes share, and the formula then overestimates two-hop success by
57 to 64 percentage points. A null-model test shows the structural counterpart:
the degree attack outperforms configuration-model theory because PrimeKG hangs
its low-degree entities, including most drugs, on its hubs. Degree correlations
explain about 40\% of that excess fragility, and structure beyond the joint
degree distribution explains about half.

\paragraph{An interpretation.}
The analysis suggests a simple way to think about redundancy. A fact is
protected not by copies but by \emph{independent} routes, and $m$ routes give
$m-1$ independent cycles. Each cycle is an additional constraint that lets a
lost edge be bypassed, much as a parity check lets an erased bit be recovered
in an error-correcting code. The analogy has a precise counterpart in physics:
for the toric code, the tolerable rate of qubit loss equals the bond percolation
threshold of the square lattice \cite{stace2009thresholds}. We use this only as
an analogy, not as a claim that knowledge graphs are codes. It does suggest the
right design question for a knowledge graph: how many independent routes
protect each fact, and how many of them share a hub.

\subsection{Limitations}
The standard 5,000-node graph is a deterministic BFS sample, not a representative random sample of the filtered PrimeKG network. At 90\% random edge removal, its giant-component fraction was 81.5\%, compared with 55.4\% for the filtered full graph. This comparison shows a substantial sampling sensitivity, but does not isolate the individual effects of graph size, composition, and degree distribution.

Our functional benchmark consists of 200 initially retrievable drug--disease pairs, balanced between single-path and multi-path cases. It therefore measures survival under a controlled task distribution, not the prevalence of retrievable relations among arbitrary pairs. The hop-limit scan uses only three random edge-removal seeds, while several other comparisons use ten; results should be interpreted with their stated seed counts, and Section~\ref{sec:prediction-uncertainty} gives long-run means from 600 fresh orders. A surviving typed association path also does not establish a clinically meaningful treatment relationship.

The targeted attacks use centrality rankings computed on the intact graph, without recomputation after removal. The retrieval model tests path availability under specified relation and hop constraints rather than a complete biomedical reasoning system. Restoration is evaluated on the BFS subgraph, and the task-aware policy uses a one-time two-hop score based on training paths missing one edge; it does not update scores during repair or consider paths missing both edges. The 140/60 training/test split separates query IDs, not entities, so endpoints may overlap across splits. These choices support controlled comparisons but limit generalization.

The theoretical results have their own limits. The percolation and $L=\infty$
formulas assume an uncorrelated, locally tree-like graph; the null-model test
shows how this assumption fails under attack, but it does not separate
clustering, disease modules and BFS sampling within the remaining half of the
gap. The Molloy--Reed threshold ($p_c\approx0.006$) is an estimate; our scan
stops at $p=0.01$, so we did not measure it. The two-hop formula is exact only
for blind damage and for routes that share no edge, which holds here because
every endpoint pair carries one relation. Per-query prediction under attack is
exploratory, because the degree features include the quantity that defines the
attack. The null models use one fixed ranking, five realisations each and ten
attempted swaps per edge.

\subsection{Conclusion}
A knowledge graph can remain almost fully connected while losing many of the
short, typed paths that its queries need. On PrimeKG, removing half of the edges
leaves 98\% of nodes connected but 43\% of the benchmark queries answerable. The gap
has a structural reason, which should carry over to other knowledge graphs even
though its size will not: the two measures are different kinds of quantity. Connectivity is a collective property with a distant percolation
threshold. Bounded-hop retrieval is a local reliability problem, governed by
how many independent routes protect each fact. A one-line formula built on that
idea predicts retrieval under random damage without fitted parameters, explains
which repairs help, and identifies exactly when it should fail: when damage
targets the hubs through which those routes pass.

A knowledge graph does not forget when it falls apart; it forgets when a fact
runs out of independent ways to be reached. Its robustness is therefore not a
property of the network alone. It depends on how many independent routes
protect each fact, how far the task is allowed to reach, and whether the damage
is blind or aimed at the hubs those routes share. A practical robustness
assessment should report task-level retrieval alongside connectivity, count
independent routes rather than paths, and state the graph, the queries, the
horizon and the damage model explicitly.
